% Chapter 20
\chapter{The Fall of Space}

\chapterepigraph{%
  Gravity is not a force pulling things together.\\
  It is the geometry of where futures are.
}{Flyxion, \textit{Axioms for a Falling Universe}}

\sidenote{RSVP accessibility\\as emergent gravity;\\ connects to\\Verlinde (2011)}

This chapter develops the most radical claim of the RSVP framework:
that space itself is a dynamical entity that ``falls'' — relaxes toward
configurations of lower constraint tension. Gravity, in this picture,
is not a fundamental force but an emergent consequence of accessibility
geometry.

\section{Space as Relaxation}

In general relativity, gravity is spacetime curvature. Matter curves
spacetime; curved spacetime deflects matter. The relationship is
geometric and exact.

In RSVP, gravity is an \emph{emergent} phenomenon: the apparent attraction
between massive bodies is the result of both bodies flowing toward the
same accessibility sink. They do not attract each other; they are both
attracted by the same accessibility gradient.

\begin{definition}{Emergent Gravitational Acceleration}{emergent-gravity}
  In the RSVP framework, the acceleration of a test mass at position $x$
  is:
  \[
    \mathbf{a}(x) = \gamma \nabla S(x)
  \]
  where $\gamma$ is the coupling constant and $\nabla S(x)$ is the
  accessibility gradient at $x$.
\end{definition}

This is formally identical to the RSVP vector field equation's
accessibility driving term (Chapter~18, Eq.~\ref{eq:v-rsvp-ch18}).
The accessibility gradient plays the role of the gravitational potential
gradient $-\nabla \Phi_\text{grav}$ in Newtonian gravity.

\section{Structure Without Expansion}

A major challenge for any non-expanding cosmology is explaining the
observed large-scale structure: why matter is concentrated in filaments,
sheets, and clusters rather than distributed uniformly.

In RSVP, structure formation is explained by accessibility landscape
topography: matter flows toward lamphrodyne regions (accessibility sinks),
creating the observed filamentary structure. The cosmic web emerges
not from gravitational instability in an expanding medium but from
accessibility gradient flow in a relaxing plenum.

\begin{definition}{Cosmic Web in RSVP}{cosmic-web-rsvp}
  The \emph{cosmic web} in the RSVP framework is the network of
  lamphrodyne regions (accessibility sinks) connected by accessibility
  ridges along which matter flows. Voids correspond to lamphron regions
  (accessibility maxima) from which matter has been driven away.
\end{definition}

\section{Constraint Descent}

The ``fall of space'' metaphor captures the RSVP picture precisely:
the universe does not expand outward; it descends along constraint
gradients. Every constraint that can be relaxed is relaxed. Every
structure that can form does form, because structure formation reduces
constraint violation at the local level (forming a stable configuration
from an unstable one).

\begin{namedthm}{The Constraint Descent Cosmology Principle}
  In the RSVP framework, cosmic evolution is constraint descent: the
  universe evolves by continuously reducing its total constraint
  violation functional $\int_M \kappa(\Phi,\mathbf{v},S)\,d^4x$.
  Gravity, structure formation, the arrow of time, and the eventual
  heat death are all consequences of this descent. Space does not
  expand; it falls.
\end{namedthm}

\section{Observable Consequences}

The RSVP framework makes specific predictions that differ from standard
cosmology:

\begin{enumerate}
  \item \textbf{Redshift-structure correlation:} redshift should correlate
        with local accessibility gradient magnitude, not just distance.
        Regions of high local structure density should show systematic
        redshift deviations.
  \item \textbf{Anisotropic redshift:} photons traveling along accessibility
        ridges should show different redshifts than photons traversing
        voids, even at the same distance.
  \item \textbf{No dark energy:} the accelerating expansion attributed
        to dark energy in standard cosmology is reinterpreted as
        increasing accessibility gradient magnitude in the late universe.
\end{enumerate}

\begin{exercise}
  In Newtonian gravity, the gravitational potential $\Phi_\text{grav}$
  satisfies $\nabla^2 \Phi_\text{grav} = 4\pi G \rho$
  (Poisson equation). In RSVP, if $\mathbf{a} = \gamma \nabla S$,
  what equation does $S$ satisfy in the Newtonian limit? Is it the
  same as Poisson's equation? What is the RSVP analog of the density $\rho$?
\end{exercise}

\begin{exercise}
  Verlinde (2011) proposed that gravity is an entropic force arising
  from the holographic principle. Compare this to the RSVP accessibility
  gradient interpretation. Where do the two proposals agree? Where do
  they differ?
\end{exercise}
